\textbf{Motivations. }
Most well-known quantum algorithms, such as Shor's algorithm~\cite{Sho97}, were
historically designed on the Quantum Random Access Machine (QRAM)
model~\cite{Kni22}. In this model, a program
interacts with a quantum memory through basic operations complying with the
laws of quantum mechanics. These operations include a fixed set of quantum
gates, chosen to be universal, as well as a probabilistic measurement of
qubits. Consequently, the control flow is purely classical, \ie depends on the
(classical) outcome of a measure. Based on  this paradigm, several high-level
quantum programming languages have been introduced, each with different
purposes and applications, from assembly code~\cite{QASM17,QASM22}, to
imperative languages~\cite{FY21}, circuit description languages~\cite{GLRSV13},
and $\lambda$-calculi~\cite{SV09}.

A relevant issue was to extend these models to programs with \emph{quantum
control}, also known as coherent control, enabling a ``program as data''
treatment for the quantum paradigm. Quantum control consists in the ability to
write superposition of programs in addition to superposition of data and
increases the expressive power of quantum programming languages. It allows the
programmer to write algorithms such as the \emph{quantum switch}, which uses
fewer resources (quantum gates) than algorithms with classical
control~\cite{CAPV13} and is physically implementable~\cite{AWHM20,PMA+15}.
Hence, quantum control provides a computational advantage over classical
control~\cite{ACB14,TCM+21,KOY24}. In the last decades, several  quantum
programming languages implementing this concept have been designed,
non-exhaustively~\cite{AG05,SVV18,DCM22,Lem24, YVC24}. To get the best of both worlds,
a natural next step was the development of \emph{hybrid languages}, \ie
languages that allow classical and quantum flow/data to be combined.
The hybrid paradigm is conveying considerable practical interest: for example,
it is used by quantum variational algorithms, a class of quantum algorithms
leveraging both classical and quantum computing resources to find approximate
solutions to optimization problems; and the properties of hybrid languages have
also been deeply studied, non-exhaustively~\cite{Qunity23,Yin24,DLPZ25,BPP25}.

Since hybrid languages offer interesting prospects in terms of expressiveness
and optimal resource consumption, a relevant and open issue concerns the
development of resource-aware (static) analyses of their programs. These static
analyses can be applied to predict the low-level resources required to execute
quantum algorithms, such as the depth or size of a quantum circuit.

\medskip
\noindent
\textbf{Contributions. }
This work solves the above issue for the first time by introducing a HYbrid
Recursive Quantum Language \hyrql{} on which a resource analysis can be
performed. \hyrql{} is a hybrid extension of the functional quantum language
$\spm$ (\emph{Symmetrical pattern-matching}) of~\cite{SVV18,Lem24}. In $\spm$,
the programmer can write down directly unitary applications, combining quantum
superposition of terms with pattern-matching. Unitarity is enforced by a
(decidable) linear typing discipline using an orthogonality predicate. \spm{}
is, however, a purely quantum language, which makes the manipulation (and,
consequently, resource analysis) of classical information awkward: they can be
neither discarded nor duplicated by linearity. The hybrid nature of \hyrql{}
relies on a typing discipline which delineates a clear separation between
quantum/linear and classical/non-linear data in the typing contexts. Similarly
to $\spm$, $\hyrql$ does not include measurement: the flow from quantum data to
classical data is handled through the use of a \textshape\ construct,  as
introduced in \cite{GLRSV13}, which returns the classic structural information
on data containing quantum information, without getting any information on the
value of the quantum states. For example, the  \textshape\ of a qubit list
provides classical (duplicable) information, such as the number of qubits. The
decision to provide $\hyrql$ with a $\spm$-architecture was motivated by two
important factors. First, in contrast with most of its competitors, $\spm$ does
not require the inclusion of an initial set of quantum gates, making its
resource analysis generic. Second, thanks to its pattern-matching design, $\spm
$ allows for the reuse of a wide variety of tools designed for resource
analysis (e.g., termination or complexity) of term rewriting systems.

Our paper contains the following main contributions:

\begin{itemize}
  \item The introduction of the hybrid language with general recursion $\hyrql$,
    as well as its operational semantics, type system, and
    standard properties: confluence (Theorem~\ref{thm:confluence}), progress
    (Theorem~\ref{lem:progress}), and subject reduction
    (Theorem~\ref{lem:subred}).
  \item A proof that the orthogonality predicate in $\hyrql$ is
    $\Pi^0_2$-complete, hence undecidable
    (Theorem~\ref{thm:orthogonality-undecidable}). Consequently, typing is not
    decidable. This is not surprising as the language encompasses general
    recursion. The decidability of typing can still be recovered, on expressive
    sub-fragments of the language, e.g., when $\hyrql$ is restricted to finite
    types and terminating programs (Proposition~\ref{prop:orthogonality-lower}).
  \item A compilation from terms (of the good type) to quantum circuits, whose
    size is bounded by the runtime-complexity of the initial term
    (Theorem~\ref{thm:circuit-bound}). In particular, terms that terminate in
    polynomial time characterize exactly the complexity class \fbqp{}
    (Theorem~\ref{thm:fbqp-sound} and~\ref{thm:fbqp-complete}). This result is
    not trivial, as lambda calculi do not comply with van Emde Boas' invariant
    thesis \cite{VEB90}, i.e., a polynomial time reduction on a lambda term
    may correspond to an exponential time reduction on a Turing machine.

  \item An overview on how complexity results can be obtained and
    automatized by translating our language into term-rewrite systems, to use
    existing techniques for runtime analysis (e.g.,
    interpretations~\cite{MP09}, recursive path orderings~\cite{Der82,Der87},
    dependency pairs~\cite{AG00}, size-change principle \cite{LJBA01}).
\end{itemize}

We illustrate the expressive power of the language through several examples:
implementa\-tion of basic quantum gates (Hadamard gate,
Example~\ref{ex:had-syntax}), higher-order (quantum switch,
Example~\ref{ex:qs-syntax}), general recursion (bounded quantum walk,
Example~\ref{ex:bqwalk}), and hybrid control flow (protocol
\emph{BB84}~\cite{BB84}, Example~\ref{ex:hybrid-function}).


\newcommand{\dlpz}{{\texttt{DLPZ}}}
\newcommand{\foq}{{\texttt{FOQ}}}
\newcommand{\qugcl}{{\texttt{QuGCL}}}
\newcommand{\qunity}{{\texttt{Qunity}}}
\newcommand{\rqc}{{\texttt{RQC}}}

\begingroup
\newcommand{\cmark}{\ding{51}}
\newcommand{\xmark}{\ding{55}}
\renewcommand{\arraystretch}{1.2}
\setlength\tabcolsep{3pt}
\begin{table}[!h]
  \begin{center}
      \begin{tabular}{c | c | c | c | c | c | c }
        & \hyrql{} &\foq{} \cite{HPS23} & $\rqc^{++}$ \cite{YZ24} &
        $\mathtt{Silq}$ \cite{Silq20} & \qunity{} \cite{Qunity23} & \dlpz{}
        \cite{DLPZ25} \\ \hline 
        Resource-aware & \cmark & \cmark & \xmark & \xmark & \xmark & \xmark \\
        \hline
        General recursion & \cmark & \xmark & \cmark & \xmark & \xmark & \xmark
        \\
        Higher-order & \cmark & \xmark & \xmark & \xmark & \xmark & \xmark
      \end{tabular}
  \end{center}
  \caption{Comparison table between hybrid quantum languages }
  \label{tab:related}
  \vspace*{-0.5cm}
\end{table}
\endgroup

\medskip
\noindent
\textbf{Related works. }
In Table~\ref{tab:related}, we provide a non-exhaustive comparison between the
main families of hybrid quantum languages which treat classical data as
first-class citizen.
Towards that end, we consider the three following criteria.
\begin{itemize} 
  \item \textit{Resource-aware} highlights if the language is designed towards
    resource analysis, that is, allows to characterize well-known complex
    classes of quantum computational models. This is the case of
    $\foq$~\cite{HPS23,HPS25}, which characterizes, under restrictions, quantum
    polynomial time and quantum polylogarithmic time~\cite{FHPS25}.
    However, $\hyrql$ has a greater expressive power than $\foq$  as it is not
    restricted to unitary operators. Moreover, all standard inductive datatypes
    can be expressed in $\hyrql$, whereas $\texttt{FOQ}$ is restricted to lists
    of qubits. 

  \item \textit{General recursion} specifies whether the language includes
    unbounded recursion on classical data. The ability to handle this type of recursion is
    necessary for a resource analysis to be relevant (\ie not trivial as in
    the case of strong normalizing languages). 

  \item \textit{Higher-order} denotes the ability of the language to
    feature general higher-order terms, i.e., functions taking functions as
    input. While some languages of Table~\ref{tab:related} have syntactic
    constructs to do so, most of them are restricted to first-order calls. 
\end{itemize}

As we are interested in resource analysis in the general setting,
Table~\ref{tab:related} only compares \hyrql{} with hybrid languages, rather
than with existing languages with quantum control, non-exhaustively
\cite{SVV18,DCM22,Yin24,YVC24}.
Other studies have already been carried out on resource analysis in quantum
programming languages. They use techniques (\eg type systems) developed in the
field of Implicit Computational Complexity (ICC, see~\cite{Pec20}) to
characterize quantum complexity classes~\cite{DLMZ10}, to infer the expected
cost or expected value of a quantum program~\cite{AMPPZ22}, or to infer upper
bounds on quantum resources (depth and size of circuits)~\cite{CDL25}. However,
these studies are restricted to classical control. For example, \cite{DLMZ10}
studies the use of soft linear logic~\cite{L04} on a quantum lambda
calculus~\cite{SV05}, \cite{AMPPZ22} adapts expectation transformers on a
language based on $\texttt{QPL}$~\cite{Sel04}, and~\cite{CDL25} develops a
dependent type system on a variant of the Quipper circuit description
language~\cite{GLRSV13}.
